\SourceLeaf{227}
form a spherical triangle. In the second case they lie in the same order as the points \((1),(2),(3)\), and in the third they lie in the reverse order. If we denote the angles of this triangle simply by \(L,L',L''\) and the perpendicular drawn on the spherical surface from the point \(L''\) to the side \(LL'\) by \(p\), then
\[
\sin p=\sin L\mathord{.}\sin LL''=\sin L'\mathord{.}\sin L'L'',\quad\text{and}\quad\lambda L''=90^0\mp p
\]
\noindent where the upper sign applies in the second case and the lower sign in the third. We therefore conclude that
\[
\begin{aligned}
\pm\Delta&=\sin L\mathord{.}\sin LL'\mathord{.}\sin LL''=\sin L'\mathord{.}\sin LL'\mathord{.}\sin L'L''\\
&=\sin L''\mathord{.}\sin LL''\mathord{.}\sin L'L''
\end{aligned}
\]
\noindent The first case can evidently be regarded as included in the second or third. It is easy to see that \(\pm\Delta\) gives six times the volume of the pyramid formed by the points \(L,L',L''\) and the centre of the sphere. Finally, we can readily conclude that the same expression \(\pm\tfrac{1}{6}\Delta\) generally gives the volume of any pyramid formed by the origin of coordinates and the points with coordinates \(x,y,z;\ x',y',z';\ x'',y'',z''\).

\Article{3}
\indent A curved surface is said to have continuous curvature at a point \(A\) on it if the directions of all straight lines drawn from \(A\) to all points of the surface infinitely close to \(A\) depart by an infinitely small amount from one and the same plane through \(A\). This plane is said to \textit{touch} the curved surface at \(A\). If this condition cannot be met at a point, continuity of curvature is interrupted there, as happens, for example, at the apex of a cone. The present investigations will be confined to curved surfaces, or to parts of surfaces, where continuity of curvature is nowhere interrupted. Here we only note that at singular points, where continuity of curvature is interrupted, the methods for determining the position of the tangent plane cease to be effective and must lead to indeterminate results.

\Article{4}
\indent The orientation of the tangent plane is most easily found from the direction of the line perpendicular to it at \(A\). This line is also called the normal to the curved surface itself. We shall represent the direction of this normal by the point \(L\) on the auxiliary sphere and set
